\chapter*{Abstract}
\addcontentsline{toc}{chapter}{Abstract}
\markboth{Abstract}{Abstract}

A large-language-model program-repair loop alternates between a proposer, which
draws a candidate patch, and an oracle, which refutes it by finding a test case
on which the patch is wrong. The proposer is the half that gets the attention
and the oracle is the half that costs. On the corpus studied here the oracle's
own reference program takes a median of 0.076\,s of CPU per test
case but a 99th percentile of 4.48\,s --- a spread of
$58.7\times$ inside the same pool --- and a loop that walks that pool in the
order the benchmark ships spends 2.18\,s of CPU per round, on every
round, for twenty rounds.

Existing memories for repair loops do not touch that cost. They are keyed by a
failure \emph{class}: fault-recorded testing prioritisation ranks a case by how
many patches it has already falsified, and the predecessor of this work,
\cegmem, stored each refuted patch under a type formed from its edit location
and the property it violated. A class-keyed memory can say which patches not to
try again --- \cegmem\ cut oracle calls per episode from $11.11$ to $1.57$,
a factor of $7.07$ --- but it cannot say which test to run first, and it cannot
say which test is cheap.

This proposal develops \caster, a repair loop whose memory is keyed by
\textbf{test case} rather than by failure class. Each case carries an online
Beta-posterior kill rate $K$ and an online mean CPU cost $\hat{m}$, and probes
are scheduled by the index \km: refute early, and refute cheaply. Keying on the
case is also what makes the second half possible, because a case is an
\emph{observation} and only an observation can be phrased to the proposer as a
positive requirement rather than a prohibition.

\paragraph{What this document reports as evidence.} The preliminary validation
is a pilot on the frozen data of the \cegmem\ study: $68$ ConDefects tasks,
$2{,}523$ test cases and $4{,}670$ stored patches for which a complete
patch~$\times$~case verdict matrix, a measured per-case reference CPU cost, and
a per-case coverage trace all exist. Two experiments run on it. A replay
simulation reproduces an arbitrary probe order exactly from the verdict matrix;
a second study re-executes the loop for real over $6{,}800$ rounds per policy
and measures CPU with \code{getrusage}. The measured result is that \km\ costs
\textbf{76.2\,\%} less oracle CPU than the shipped pool order
(95\,\% CI $[-84.9, -61.2]$, $p = 1.7\times10^{-11}$, cheaper on $63$ of $68$
tasks) and \textbf{15.7\,\%} less than the published
fault-recorded baseline (95\,\% CI $[-28.9, -5.1]$, $p = 0.011$, $42$ of $68$)
--- while paying 23.3\,\% \emph{more} probes than that baseline.
Executions and CPU therefore disagree in sign on this corpus, which matters
because the test-prioritisation literature deliberately reports executions.

\paragraph{What it does not report.} The \caster\ implementation itself has
been built and unit-tested but its experimental grid has not been run. The
twenty-two arms, the steered prompt contract that the test-case key is supposed
to enable, the duplicate-rate and token-cost measurements, and the full
fault corpus are Phase~2 work and are stated as such throughout. No number in
Chapter~\ref{ch:evaluation} is a result; every number there is a target.
